\chapter{Yang--Baxter Consistency and Internal Symmetry}
\label{ch:integrable-yang-baxter-internal-symmetry}

\ConventionsLorentzian

The two-body bootstrap datum of
Chapter~\ref{ch:integrable-scattering-analyticity-bootstrap} becomes a
many-particle scattering theory only if different decompositions of the same
multi-particle process agree.  This chapter formulates that requirement as an
identity of meromorphic maps on species spaces and explains how internal
symmetry reduces the identity to scalar functions on irreducible channels.

\section{Species Spaces and Two-Body Operators}

Let \(\mathcal I\) be a set of stable particle species.  To each
\(a\in\mathcal I\) assign a finite-dimensional complex vector space \(V_a\),
called the species space.  The one-particle momentum variable of a particle of
mass \(m_a\) is the rapidity \(\theta\), with
\[
  p_a(\theta)=m_a(\cosh\theta,\sinh\theta).
\]
For ordered species \(a,b\), a two-body scattering operator is a meromorphic
map
\[
  S_{ab}(\theta):V_a\otimes V_b\longrightarrow V_b\otimes V_a,
  \qquad \theta=\theta_a-\theta_b .
\]
The codomain order is part of the definition: after scattering in one spatial
dimension the outgoing ordered particles have exchanged their positions.

It is often convenient to compose \(S_{ab}\) with the ordinary flip
\(P_{ab}:V_b\otimes V_a\to V_a\otimes V_b\) and define
\[
  R_{ab}(\theta)=P_{ab}S_{ab}(\theta):
  V_a\otimes V_b\longrightarrow V_a\otimes V_b .
\]
The \(S\)-operator records physical ordering, while \(R\) records an
endomorphism on a fixed tensor product.  All equations below are stated in the
\(R\)-matrix convention because the domains and codomains are then identical.

\begin{definition}[Regular two-body \(R\)-datum]
\label{def:regular-two-body-r-datum}
A regular two-body \(R\)-datum consists of species spaces \(V_a\) and
meromorphic endomorphisms \(R_{ab}(\theta)\in
\operatorname{End}(V_a\otimes V_b)\) such that, on every common domain of
holomorphy,
\[
  R_{ab}(\theta)R_{ab}(-\theta)=\id_{V_a\otimes V_b}
\]
and the crossing and residue data of
Definition~\ref{def:factorized-scattering-bootstrap-datum} are satisfied
after undoing the flip.
\end{definition}

\section{Derivation of the Yang--Baxter Equation}

Consider three particles with rapidities
\(\theta_1>\theta_2>\theta_3\) and species spaces
\(V_{a_1},V_{a_2},V_{a_3}\).  Factorization asserts that the full
three-particle scattering is a product of two-body scatterings.  There are two
reduced decompositions of the permutation reversing the order of the three
particles:
\[
  (12)(23)(12)=(23)(12)(23).
\]
The equality of the resulting linear maps on
\[
  V_{a_1}\otimes V_{a_2}\otimes V_{a_3}
\]
is therefore the condition that factorized scattering be independent of the
chosen reduced decomposition.

Write \(\theta_{ij}=\theta_i-\theta_j\).  Let \(R_{ij}\) denote the
endomorphism \(R\) acting on the \(i\)-th and \(j\)-th tensor factors and as
the identity on the remaining factor.  Equality of the two decompositions is
\[
  R_{12}(\theta_{12})R_{13}(\theta_{13})R_{23}(\theta_{23})
  =
  R_{23}(\theta_{23})R_{13}(\theta_{13})R_{12}(\theta_{12}).
\]
Since \(\theta_{13}=\theta_{12}+\theta_{23}\), this is the spectral-parameter
Yang--Baxter equation
\[
  R_{12}(\theta)R_{13}(\theta+\phi)R_{23}(\phi)
  =
  R_{23}(\phi)R_{13}(\theta+\phi)R_{12}(\theta).
\]

\begin{proposition}[Three-particle consistency]
\label{prop:yb-three-particle-consistency}
For a regular two-body \(R\)-datum, the factorized three-particle scattering
operator is independent of the reduced decomposition of the three-particle
permutation if and only if the Yang--Baxter equation holds as a meromorphic
identity on every triple species space.
\end{proposition}

\begin{proof}
Each adjacent transposition is represented by the corresponding two-body
operator with the rapidity difference of the two lines at the moment they
cross.  In the \(R\)-matrix convention this representation acts on a fixed
tensor product.  The two reduced decompositions of the longest element of the
three-particle permutation group give exactly the two sides of the displayed
Yang--Baxter equation.  Thus independence of the decomposition gives the
equation.  Conversely, if the equation holds, the two possible three-particle
factorizations define the same meromorphic map.  The remaining relations of
the permutation group are the inverse relation, supplied by unitarity, and
commutation of disjoint transpositions, which is automatic because the
operators act on disjoint tensor factors.
\end{proof}

\section{Internal Symmetry and Projector Decomposition}

Let \(G\) be a group or Lie algebra acting on the species spaces by
representations \(\rho_a:G\to GL(V_a)\).  Internal-symmetry invariance of the
two-body scattering datum means
\[
  R_{ab}(\theta)(\rho_a(g)\otimes\rho_b(g))
  =
  (\rho_a(g)\otimes\rho_b(g))R_{ab}(\theta),
  \qquad g\in G .
\]
Equivalently, \(R_{ab}(\theta)\) is an element of
\[
  \operatorname{End}_G(V_a\otimes V_b),
\]
the algebra of \(G\)-intertwiners.  If
\[
  V_a\otimes V_b\simeq \bigoplus_{\lambda}
  M_{ab}^{\lambda}\otimes W_\lambda
\]
is the decomposition into irreducible \(G\)-modules \(W_\lambda\) with
multiplicity spaces \(M_{ab}^{\lambda}\), then Schur's lemma gives
\[
  \operatorname{End}_G(V_a\otimes V_b)
  \simeq
  \bigoplus_{\lambda}\operatorname{End}(M_{ab}^{\lambda}).
\]
If every multiplicity space is one-dimensional, the most general invariant
\(R\)-matrix has the form
\[
  R_{ab}(\theta)=\sum_{\lambda} r_{ab}^{\lambda}(\theta)P_{\lambda},
\]
where \(P_\lambda\) is the projector onto the irreducible channel
\(\lambda\), and \(r_{ab}^{\lambda}\) is a scalar meromorphic function.

\begin{example}[Vector representation of \(O(N)\)]
\label{ex:on-projector-channels}
Let \(V=\C^N\) be the vector representation of \(O(N)\).  The tensor product
decomposes as
\[
  V\otimes V
  =
  \mathbf 1\oplus \Lambda^2 V\oplus
  \operatorname{Sym}^2_0 V .
\]
An \(O(N)\)-invariant two-body \(R\)-matrix without multiplicity mixing is
therefore
\[
  R(\theta)=r_0(\theta)P_0+r_A(\theta)P_A+r_T(\theta)P_T,
\]
with projectors onto the singlet, antisymmetric, and traceless-symmetric
channels.  The Yang--Baxter equation becomes a coupled system of functional
equations for \(r_0,r_A,r_T\), together with the analytic constraints of the
bootstrap datum.
\end{example}

\section{What the Equation Certifies}

The Yang--Baxter equation certifies compatibility of factorized
multi-particle on-shell scattering.  It does not by itself construct local
operator algebras, form factors, or a Hilbert-space QFT.  The next chapter
adds the form-factor problem, where local operators are represented by
meromorphic functions satisfying exchange, crossing, residue, and growth
conditions.

\begin{openproblem}[From \(R\)-matrices to local fields]
\label{op:r-matrix-to-local-fields}
Given a meromorphic \(R\)-datum satisfying unitarity, crossing, the
Yang--Baxter equation, residue positivity, and stated growth bounds, construct
or obstruct a local QFT realization.  The construction must identify a Hilbert
space, wedge-local or local algebras, form factors of a dense set of local
operators, and the topology in which locality and clustering are verified.
\end{openproblem}
